\begin{examplebox}
	\textbf{\textcolor{CExample}{例 1（基础）}}

	\textbf{\textcolor{CExample}{题目：}}
	如图，直二面角 $\alpha-l-\beta$，点 $A\in\alpha$，$AC\perp l$ 于 $C$。求证：$AC\perp\beta$。

	\textbf{\textcolor{CExample}{解题步骤：}}
	\begin{description}[style=nextline, leftmargin=2.8em, labelsep=0.5em, itemsep=0.45em]
		\item[\textbf{第一步（认定面面垂直）：}] 已知 $\alpha\perp\beta$，交线为 $l$。
			\[
				\alpha\perp\beta,\ \alpha\cap\beta=l.
			\]
			\par\textbf{理由：} 题设条件直接给出。

		\item[\textbf{第二步（认定线面条件）：}] $AC\subset\alpha$，且 $AC\perp l$。
			\[
				AC\subset\alpha,\quad AC\perp l.
			\]
			\par\textbf{理由：} 点 $A\in\alpha$，$AC\perp l$ 于 $C$，$C\in l$，所以 $AC$ 在 $\alpha$ 内且垂直于 $l$。

		\item[\textbf{第三步（应用定理）：}] 由面面垂直性质定理，若两平面垂直，在一个平面内垂直于交线的直线垂直于另一个平面，因此 $AC\perp\beta$。
			\[
				\therefore AC\perp\beta.
			\]
			\par\textbf{理由：} 满足定理的条件，直接得到结论。
	\end{description}

	\textbf{\textcolor{CExample}{关键结论：}}
	面面垂直性质定理可直接用于证明线面垂直。

	\medskip
	\textbf{\textcolor{CExample}{例 2（稍变形）}}

	\textbf{\textcolor{CExample}{题目：}}
	在四棱锥 $P-ABCD$ 中，底面 $ABCD$ 是正方形，侧面 $PAD\perp$ 底面 $ABCD$，且 $PA=PD$，$E$ 为 $AD$ 中点。求证：$PE\perp$ 底面 $ABCD$。

	\textbf{\textcolor{CExample}{解题步骤：}}
	\begin{description}[style=nextline, leftmargin=2.8em, labelsep=0.5em, itemsep=0.45em]
		\item[\textbf{第一步（证等腰中线垂直）：}] 因为 $PA=PD$，$E$ 为 $AD$ 中点，所以 $PE\perp AD$。
			\[
				PA=PD,\ AE=ED\ \Rightarrow\ PE\perp AD.
			\]
			\par\textbf{理由：} 等腰三角形性质（三线合一）。

		\item[\textbf{第二步（设定符号与位置）：}] 记侧面 $PAD$ 为 $\alpha$，底面 $ABCD$ 为 $\beta$，交线 $AD$ 为 $l$。则 $PE\subset\alpha$，$l=\alpha\cap\beta$，且 $PE\perp l$。
			\[
				PE\subset\alpha,\ l=\alpha\cap\beta,\ PE\perp l.
			\]
			\par\textbf{理由：} 第一步结论及已知。

		\item[\textbf{第三步（确认面面垂直）：}] 已知 $\alpha\perp\beta$（即侧面 $PAD\perp$ 底面 $ABCD$）。
			\[
				\alpha\perp\beta.
			\]
			\par\textbf{理由：} 题设条件直接给出。

		\item[\textbf{第四步（应用定理）：}] 由面面垂直性质定理，$PE\perp\beta$，即 $PE\perp$底面 $ABCD$。
			\[
				\therefore\ PE\perp\beta.
			\]
			\par\textbf{理由：} 定理：若两平面垂直，在一平面内且垂直于交线的直线必垂直于另一平面。
	\end{description}

	\textbf{\textcolor{CExample}{关键结论：}}
	在面面垂直条件下，先证线线垂直再套定理，是常见思路。

	\medskip
	\textbf{\textcolor{CExample}{例 3（综合应用）}}

	\textbf{\textcolor{CExample}{题目：}}
	在等腰梯形 $ABCD$ 中，$AD\parallel BC$，$AD=1$，$BC=3$，$\angle ABC=60^\circ$。将 $\triangle ABD$ 沿 $BD$ 折起，使平面 $ABD\perp$ 平面 $BCD$。在折起后的图形中，过点 $A$ 作 $AO\perp BD$ 于点 $O$。求证：$AO\perp$ 平面 $BCD$，并求点 $A$ 到平面 $BCD$ 的距离。

	\textbf{\textcolor{CExample}{解题步骤：}}
	\begin{description}[style=nextline, leftmargin=2.8em, labelsep=0.5em, itemsep=0.45em]
		\item[\textbf{第一步（作辅助垂线）：}] 在平面 $ABD$ 内，过 $A$ 作 $AO\perp BD$ 于 $O$。
			\[
				AO\perp BD.
			\]
			\par\textbf{理由：} 过一点可作直线的垂线，垂足为 $O$。

		\item[\textbf{第二步（应用定理得线面垂直）：}] 记平面 $ABD$ 为 $\alpha$，平面 $BCD$ 为 $\beta$，交线 $BD$ 为 $l$。则 $\alpha\perp\beta$，$l=\alpha\cap\beta$，$AO\subset\alpha$，$AO\perp l$。由面面垂直性质定理，得 $AO\perp\beta$，即 $AO\perp$平面 $BCD$。
			\[
				AO\perp\beta.
			\]
			\par\textbf{理由：} 面面垂直性质定理。

		\item[\textbf{第三步（求 $AO$ 长度）：}] 在折叠前的等腰梯形中，过 $A$ 作 $AE\perp BC$ 于 $E$，则 $BE=\frac{BC-AD}{2}$，故 $BE=1$。由 $\angle ABC=60^\circ$，得 $AB=2$，$AE=\sqrt{3}$。在 $\triangle ABD$ 中，$AD=1$，$AB=2$，$\angle BAD=120^\circ$，由余弦定理：
			\[
				\begin{aligned}
					BD^2 &= AD^2+AB^2-2\cdot AD\cdot AB\cdot\cos120^\circ\\
					&= 1+4-2\times1\times2\times\left(-\frac12\right)\\
					&= 7,
			\end{aligned}
			\]
			所以 $BD=\sqrt{7}$。
			又 $\triangle ABD$ 的面积
			\[
				\begin{aligned}
					S &= \frac12 AD\cdot AB\cdot\sin120^\circ\\
					&= \frac12\times1\times2\times\frac{\sqrt3}{2}\\
					&= \frac{\sqrt3}{2}.
			\end{aligned}
			\]
			由等面积法 $S=\frac12 BD\cdot AO$，得
			\[
				\begin{aligned}
					AO &= \frac{2S}{BD}\\
					&= \frac{\sqrt3}{\sqrt7}\\
					&= \frac{\sqrt{21}}{7}.
			\end{aligned}
			\]
			\par\textbf{理由：} 利用折叠前平面图形计算边长，再通过等面积法求高。
	\end{description}

	\textbf{\textcolor{CExample}{关键结论：}}
	在折叠问题中，利用面面垂直性质作平面的垂线，再结合平面几何计算长度，是求点到平面距离的常用方法。
\end{examplebox}
